\chapter{Introduction}
\label{ch:intro}

\section{Motivation}

Large language models (LLMs) demonstrate high accuracy by storing factual knowledge in
their parameters, yet they remain prone to hallucinations, and their ability to access and
update knowledge is limited~\cite{lewis2020rag}. Retrieval-augmented generation (RAG)
addresses these shortcomings by combining the parametric language generator of an LLM with
a non-parametric external document index, from which relevant information is retrieved to
generate more accurate answers~\cite{lewis2020rag}. Within the RAG pipeline, document
chunking is the pre-processing step that splits the text into units and thereby defines the
search granularity of the vector index~\cite{smigielski2026chunking}.

Many chunking variants exist for this step, but they differ in accuracy and are not
suitable for every document type or domain. Recent research introduced a newer mechanism,
LLM-based chunking, which dynamically identifies narrative topic shifts to improve the
quality of each vector embedding~\cite{duarte2024lumberchunker}. However, LLM-driven
segmentation introduces a considerable computational latency compared to the existing
baselines, as the experiments with LumberChunker showed~\cite{duarte2024lumberchunker}.

Furthermore, empirical evidence suggests that on real-world documents, sentence-level
positional proximity already provides a strong natural heuristic for semantic shifts, so
that the computational overhead is not justified by the gain~\cite{qu-etal-2025-semantic}.
This motivates an evaluation of whether LLM-based chunking brings a real advantage over
length-matched baselines, or whether preserving the integrity of sentences is already
sufficient for effective retrieval.

\section{Problem statement and research question}
\label{sec:problem}

The research question is: \emph{Does a local language model place better chunk boundaries
than rule-based chunking techniques, and under which conditions?} Secondary questions
concern the effect of the optional features of the library: heading detection, the
low-information filter, topic enrichment and the choice of the split point when a chunk
exceeds the size cap.
